\chapter{Data Acquisition and Feature Engineering}

A robust machine learning pipeline for polymer-molecule interaction prediction requires both a reliable dataset and informative numerical representations of chemical structures. This chapter details the methodology employed to collect, augment, and transform experimental data into machine-readable feature vectors suitable for downstream modeling tasks.

\section{Literature Scraping for Polymer-Molecule Data}
\label{sec:literature_scraping}

The primary bottleneck in developing predictive models for polymer-based wastewater remediation is the scarcity of standardized experimental datasets. Unlike general molecular property databases, there exists no comprehensive public repository documenting the adsorption capacities of specific polymer-molecule pairs under controlled environmental conditions. To address this limitation, a systematic literature review was conducted to extract experimental data from peer-reviewed publications.

The data collection effort focused on identifying studies that reported quantitative measurements of pharmaceutical adsorption onto polymer materials in aqueous solutions. From each relevant publication, three key variables were extracted for each experimental condition: the pH of the solution, the initial concentration of the pharmaceutical pollutant measured in milligrams per liter (mg/L), and the resulting adsorption capacity expressed in milligrams per gram (mg/g). These values respectively encode the environmental conditions of the experiment, the thermodynamic driving force for adsorption, and the efficiency metric that characterizes the polymer's effectiveness.

The scraped data were organized into a structured CSV (Comma-Separated Values) format, where each row represents a unique experimental observation identified by the polymer used, the target pharmaceutical molecule, and the specific experimental conditions. Importantly, the original source of each data point was preserved, enabling traceability back to the primary literature. This dataset, referred to as the Polymer Drug Concentration Capacity (PDCC) dataset, serves as the foundational resource for all subsequent machine learning experiments in this thesis.

The challenges encountered during data collection are noteworthy. Experimental results are often reported in non-standardized formats, with adsorption capacities expressed using varying units or normalization schemes. Furthermore, many studies report only a limited number of data points for highly specific polymer-molecule pairs under fixed conditions, making cross-study comparisons difficult. These factors underscore the need for careful data curation and standardization before the data can be utilized for machine learning applications.

\section{Overcoming Data Scarcity: Interpolation and Augmentation}
\label{sec:data_augmentation}

The PDCC dataset, while carefully curated, remains comparatively small relative to typical machine learning benchmarks. Data scarcity poses a fundamental challenge: models trained on limited data risk overfitting to the specific experimental conditions represented in the training set, failing to generalize to novel polymer-molecule combinations. To mitigate this issue, a data augmentation strategy based on interpolation was developed and systematically evaluated.

The rationale for interpolation is grounded in the physicochemical nature of the adsorption process. When a polymer is introduced into contaminated water, the mass of pollutant absorbed per unit mass of polymer (capacity) typically exhibits a monotonic relationship with the initial pollutant concentration, approaching a saturation plateau at sufficiently high concentrations. This behavior is consistent with adsorption isotherm models such as the Langmuir or Freundlich models. Consequently, the capacity values at intermediate concentrations can be reasonably approximated by interpolation between observed data points.

The augmentation pipeline implemented in this work follows a two-stage approach. First, linear interpolation is performed between pairs of adjacent concentration-capacity data points belonging to the same polymer-molecule pair. This process generates intermediate synthetic data points that represent plausible adsorption behaviors at concentrations not explicitly measured in the original experiments. The number of interpolated points between each pair of observed values is configurable, allowing control over the degree of dataset expansion.

Second, origin points are added to each polymer-molecule group. Specifically, for each unique combination of polymer and target molecule (and accounting for pH when applicable), a synthetic observation at concentration equal to zero with capacity equal to zero is appended. This addition reflects the physically intuitive boundary condition that no pollutant can be absorbed when none is present in the solution, and it anchors the interpolated curves at the origin.

Four distinct augmentation strategies were implemented and tested: pure interpolation, interpolation followed by origin point addition, origin point addition followed by interpolation, and origin point addition only. Each method produces a different augmented dataset, and the optimal strategy was determined empirically through the cross-validation experiments described in Chapter 6.

\section{Feature Extraction using RDKit, Polymetrix, and tblite}
\label{sec:feature_extraction}

Raw SMILES and P-SMILES strings, while sufficient for generative modeling, must be converted into numerical feature vectors before they can be used as inputs to the Multi-Layer Perceptron. This section describes the feature extraction pipeline that transforms chemical line notations into chemically meaningful descriptors.

The feature extraction workflow leverages three established computational chemistry libraries: RDKit for molecular property calculation and canonicalization, Polymetrix for advanced polymer-specific descriptors, and tblite for quantum chemical calculations. Each library contributes complementary information about the molecular structure.

For both polymers and small molecules, the following descriptors are computed. The partition coefficient (logP) estimates the lipophilicity of the compound, which is critical for understanding how the molecule partitions between the aqueous phase and the polymer matrix. The distribution coefficient (logD) extends logP by accounting for ionization effects at a given pH, making it particularly relevant for wastewater applications where pH varies. HOMO-LUMO energies in electron volts (eV) characterize the electronic structure of the molecule, specifically the energy of the highest occupied molecular orbital (HOMO) and the lowest unoccupied molecular orbital (LUMO). These values are central to predicting intermolecular interactions through Frontier Molecular Orbital theory. The net charge at a specified pH represents the ionization state of the molecule under environmental conditions, influencing electrostatic attraction or repulsion with the polymer surface.

Additionally, Morgan fingerprints are computed to encode the topological structure of the molecule as a fixed-length binary vector. These fingerprints capture local chemical environments and enable the model to recognize structural similarities between molecules.

The Polymetrix library provides a suite of polymer-specific features that extend beyond standard molecular descriptors. These include the topological polar surface area (TPSA), which quantifies the total surface area of polar atoms in the molecule and correlates with hydrogen bonding capacity. The synthetic accessibility (SA) score estimates how easily a given polymer structure could be synthesized in practice, which is essential for real-world applicability. Additional Polymetrix features include counts of hydrogen bond donors and acceptors, the number of rotatable bonds, ring counts (aromatic and non-aromatic), molecular weight, and various VSA (van der Waals Surface Area) descriptors.

A notable technical challenge arises when computing features for polymers. Standard cheminformatics libraries, including RDKit, are primarily designed for discrete small molecules and may not correctly handle the non-standard atom types and connectivity patterns present in P-SMILES notation. To overcome this limitation, a capping strategy was employed: for each polymer, the repeating unit is first extracted and treated as a discrete molecular fragment. This fragment is then capped with hydrogen atoms at the attachment points, converting it into a valid molecule that can be parsed by RDKit. The features computed from this capped monomer are used as proxies for the properties of the full polymer chain. This approximation is chemically reasonable because the repeating unit determines much of the polymer's bulk chemical behavior.

The full feature extraction pipeline produces a high-dimensional numerical vector for each polymer-molecule pair, encoding structural, thermodynamic, electronic, and environmental information. These feature vectors serve as the input to the MLP described in Chapter 5, where they are scaled using a StandardScaler to ensure that all features contribute equally to the learning process.
